\bookchapter{第十七章：スマホと他のデバイスの連携}

\chapterintro{一、パソコンとの連携：スマホとパソコンはもう二つの世界ではない}

\booksection{正直、以前最も嫌だったのはスマホで撮った写真をパソコンに置いてレタッチしたくて、ケーブルを探し、転送を待つこと。後に、実はケーブル不要で、良い連携はケーブルより10倍速いことが分かりました。}

Appleエコシステム（AirDrop + ユニバーサルクリップボード）

\begin{itemize}
  \item \textbf{：iPhoneとMacを使っているなら、AirDropは最初に覚えるべき機能。スマホで共有をタップ AirDrop Macを選択、写真が瞬時に移動し、通信量もWi-Fiも不要、BluetoothとWi-Fiダイレクトを使います。もう一つ見落とされがちな機能が「ユニバーサルクリップボード」——スマホでテキストをコピーし、Macで直接ペーストすると出てきて、まるで魔法のよう。}Android＋Windows（Link to Windows）
  \item \textbf{：Microsoftには「Link to Windows」というAppがあり、Samsungスマホに内蔵、他のAndroidスマホはアプリストアからダウンロード可能。ペアリング後、パソコンでスマホの通知確認、SMS返信、写真閲覧、さらには電話の発着信が可能。2年使っていますが、会議中にスマホを何度も見る必要がなく、パソコン画面に通知がポップアップするだけで十分。}ブランドマルチスクリーン連携
  \item \textbf{：Huaweiの「マルチスクリーン連携」とXiaomiの「Xiaomi妙享」はさらに徹底——スマホ画面を直接パソコンに投影し、マウスでスマホからパソコンへファイルをドラッグ＆ドロップでき、パソコンでDouyinを見たりWeChatに返信したりも可能。クロスデバイスのコピーペースト、ファイルドラッグもでき、体験は非常にスムーズ。}二、テレビにキャスト：家族みんなで視聴
\end{itemize}

\booksection{お正月に実家に帰った時、従妹がテレビで彼女の卒業旅行動画を流したがっていました。長く苦労しましたが、最終的に非常にシンプルだと分かりました：}

iPhone テレビ

\begin{itemize}
  \item \textbf{：右上からコントロールセンターを下ろし、二つの四角が重なったアイコン（画面ミラーリング）をタップ、テレビ名を見つけてタップすれば接続。前提はテレビとスマホが同じWi-Fiに接続し、テレビがAirPlayに対応していること（ほとんどのスマートテレビが対応）。}Android テレビ
  \item \textbf{：各社呼び方が異なり、Huaweiは「ワイヤレスキャスト」、Xiaomiは「投影」、OPPO/vivoは「マルチスクリーン連携」。通常はクイックスイッチパネルで見つけられ、見つからない場合は設定で「キャスト」を検索。}動画Appから直接キャスト
  \item \textbf{：テレビでiQiyi、Bilibili、YouTubeを見たいだけなら、さらに簡単——再生画面でWi-Fiシンボル付きの四角い小アイコンを探し、タップしてテレビを選択すると動画が移動。この時スマホで別のことができ、動画はテレビで継続再生、相互に影響しません。}キャスト失敗時の対策
  \item \textbf{：90\%は同じWi-Fiに接続していないこと。残り10\%はテレビのキャスト機能がオフで、テレビ設定で「マルチスクリーン連携」や「AirPlay」スイッチを探してください。}三、スマートウォッチ/バンド：手首に常駐する健康管理者
\end{itemize}

\booksection{母にバンドを買いました。最初は「運動もしないのにこれを買ってどうするの」と言っていました。2週間使って、自分が毎日5時間ちょっとしか寝ていないことに気づき、自発的に作息を調整し始めました。これがウォッチ/バンドの最も実用的な価値——目に見えない健康データを目に見える数字にする。}

通知フィルター

\begin{itemize}
  \item \textbf{：スマホがバッグやポケットで震え続ける時、手首を上げて一目見る——上司からのWeChatかPinduoduoの販促か、1秒で判断、スマホを取り出す価値のないメッセージはそのままに。}健康データ同期
  \item \textbf{：ウォッチは24時間心拍数、血中酸素、睡眠、歩数をモニタリング、これらのデータは自動的にスマホの健康Appに集約。長期的に見れば自分のトレンドが分かる——今月は歩くのが少なすぎないか、今週の睡眠の質は良くなったか。}交通カード/キーカード同期
  \item \textbf{：スマホで登録したバスカードやキーカードは、ほとんどのブランドが自動的にウォッチに同期。地下鉄でスマホを取り出す必要がなく、手首を改札機にかざすだけで通過、本当に「未来は既に来ている」と感じます。}スマホ検索の神器
  \item \textbf{：ウォッチには通常「スマホを探す」機能があり、タップするとスマホが鳴ります。私がこの機能を使う頻度は予想をはるかに超える——スマホをどこに置いたかすぐ忘れる。}四、スマートホーム：スマホは家の「メインスイッチ」
\end{itemize}

\booksection{スマートホームは難しそうに聞こえますが、要するに3つのこと：スマホでリモート制御、自動化ルールの設定、音声操作。}

私自身のいくつかの実際の使い方をお話しすると、これがいかに実用的か分かります：

エアコン

\begin{itemize}
  \item \textbf{：夏の帰宅途中にスマホで事前にエアコンをオン、ドアを開けると涼しい。100元余りでエアコンパートナーコンセントを買えば実現でき、エアコンを買い替える必要なし。}照明
  \item \textbf{：就寝前にベッドに横になり、起きるのが面倒ならスマホでタップまたは「Hey Siri/小愛同学、電気を消して」と声かけ。}ロボット掃除機
  \item \textbf{：外出前にスマホで「掃除開始」をタップ、帰宅すると床が綺麗に。毎日定時掃除も設定可能。}カメラ
  \item \textbf{：家に老人、子供、ペットがいれば、スマホでいつでも家庭の状況を確認。}どのエコシステムを選ぶか？一言でまとめると：
\end{itemize}

iPhone使用 Apple HomeKit対応デバイスを優先（少し高いがプライバシー保護が良い）。

\begin{itemize}
  \item デバイス多数でコスパ重視 Xiaomi米家エコシステムが最も充実、照明、コンセント、ロボット掃除機など何でも揃う。
  \item Huaweiスマホ使用 HarmonyOS Connectを優先、Huaweiデバイスとの体験が最も良い。
  \item 五、Bluetoothイヤホンとスピーカー：ワイヤレスの自由
\end{itemize}

\booksection{有線イヤホンからBluetoothイヤホンへの変更は、私にとってスマホ体験のアップグレードで最も顕著なものでした。もうバッグから絡まった糸の塊を取り出すことも、ドアノブにケーブルが引っかかることもありません。}

ペアリングは簡単

\begin{itemize}
  \item \textbf{：イヤホンは初回電源オンで通常自動的にペアリングモード（点滅）に入り、スマホのBluetoothをオンにすれば検索でき、タップするだけで接続。以降は取り出すたびに自動接続、何もする必要なし。}同ブランドの相乗効果
  \item \textbf{：AirPodsとiPhoneの組み合わせはバッテリー残量が自動ポップアップし、複数のAppleデバイス間でシームレスに切り替え可能。Huawei FreeBudsとHuaweiスマホ、XiaomiイヤホンとXiaomiスマホにも同様のポップアップと最適化があります。}Bluetoothスピーカー
  \item \textbf{：家庭集会、アウトドアキャンプの時、一つのBluetoothスピーカーで雰囲気が倍増。接続後、スマホがあなたのDJブースに。}遅延に注意
  \item \textbf{：ゲーム時にBluetoothイヤホンを使うと、わずかな遅延（銃声がしてから0.5秒後に音が聞こえる）が生じる可能性。遅延に敏感な場合は、有線イヤホンを使うか、「低遅延モード」対応のBluetoothイヤホンを購入。}六、カーシステム：運転中の安全なパートナー
\end{itemize}

\begin{tcbinfo}{UWB 超広帯域技術——Bluetooth より高精度な「空間感覚」}
Bluetooth のほか、新型スマホには UWB（Ultra-Wideband、超広帯域）チップも内蔵されています。その能力は\textbf{高精度な測距と空間定位}——他のデバイスがどれくらい離れているか、どの方向にあるかを計算でき、精度はセンチメートル級に達します。Bluetooth は「近くにいる」しか分かりませんが、UWB は「左前方1.2m」まで分かります。主な用途：

\begin{itemize}
  \item \textbf{AirTag/検索の高精度化}: iPhone が AirTag を検索する際に「7m先、左へ」の矢印ナビを表示するのは UWB のおかげです。通常の Bluetooth は「範囲内にいるかいないか」しか分かりません。
  \item \textbf{デジタルカーキー}: UWB 搭載の iPhone と Android フラッグシップは、カーキーとして車に近づくと自動解除、離れると自動ロックでき、運転席と助手席を区別し、リレー攻撃も防げます。
  \item \textbf{屋内ナビゲーション}: ショッピングモール、空港で UWB ビーコンによる「通路レベル」のナビゲーションが可能で、GPS よりはるかに正確です。
  \item \textbf{サポートする機種}: iPhone 11 以降の機種に UWB チップ（U1/U2）搭載。Android 陣営では Xiaomi、Samsung、vivo、OPPO の一部フラッグシップにも搭載済み。
\end{itemize}

\end{tcbinfo}

\begin{tcbinfo}{LE Audio と Auracast——Bluetooth オーディオの次世代標準}
20年使われてきた従来の Bluetooth オーディオ（A2DP）は\textbf{LE Audio}（低消費電力オーディオ）に置き換えられつつあります。これは Bluetooth 5.2 以降の次世代オーディオ標準です：

\begin{itemize}
  \item \textbf{より省電力、より高音質}: 全新の LC3 コーデックを採用し、従来の SBC より高音質·低消費電力で、イヤホンのバッテリー持ちが顕著に延びます。
  \item \textbf{デュアルオーディオパス}: 左右のイヤーピースが独立して伝送し、片耳モードへの切り替えでも切断しません。
  \item \textbf{Auracast ブロードキャストオーディオ}: LE Audio の目玉機能。1台のスマホ/TVから同時に\textbf{無数の}イヤホンにオーディオをブロードキャスト——空港ロビー、博物館、ジム、教室で1つの送信源から、複数の人が自分のイヤホンで同じ音声を聴けます。
  \item \textbf{現状}: 2024-2026年に新発売のフラッグシップイヤホンとスマホは LE Audio を広くサポート。旧デバイスはハードウェアアップグレードが必要でソフトウェアでは解放できません。新イヤホン購入時に LE Audio/Auracast 対応表示を確認しましょう。
\end{itemize}

\end{tcbinfo}

\booksection{運転中のスマホ操作は違法で、命に関わります。しかしナビ、電話、音楽が必要な時もあります。この時にCarPlay（Apple）とAndroid Auto（Android）が役立ちます。}

接続方法

\begin{itemize}
  \item \textbf{：車が対応していれば（近年のほとんどの新車が対応）、ケーブルでスマホを車のUSBポートに接続、カーディスプレイが自動的にCarPlayまたはAndroid Auto画面に切り替わります。一部の新車はワイヤレス接続に対応、さらに便利。}できること
  \item \textbf{：ナビ（カーナビよりはるかに使いやすい）、電話（ハンドルで応答）、音楽（スマホの音楽Appを直接車に投影）、音声WeChat（Siriや小愛がメッセージの読み上げと返信をサポート）。}古い車でも利用可能
  \item \textbf{：車が古く対応していない場合、数百元で外部CarPlayディスプレイを買い、ダッシュボードに貼り付けば、体験はほぼ同じ。}七、連携中のセキュリティとプライバシー
\end{itemize}

\booksection{デバイスが多く連携されるほど、プライバシー漏洩のリスク面が拡大します。いくつか注意点：}

公共の場でキャストする前、スマホの通知バーにプライベートメッセージがないか確認、家族に見せるべきでない内容を見せないように。

\begin{itemize}
  \item デバイスを他人に譲渡または売却する前、連携解除を——スマホ設定でそのデバイスを見つけ、「削除」または「ペアリング解除」を選択。
  \item スマートホームカメラは、ブランド品で暗号化対応のものを購入、数十元の無名品は避ける。カメラがハッカーに生配信されたニュースは多すぎます。
  \item やってみよう：スマホ画面をテレビにキャストし、スマートシーンを一つ設定
\end{itemize}

\begin{tcbpractice}{スマホとテレビが同じWi-Fiに接続されていることを確認。}
\begin{enumerate}
  \item テレビで「画面ミラーリング」または「AirPlay」機能を見つけてオン。
  \item スマホで：iPhoneはコントロールセンターから「画面ミラーリング」をタップ。Androidは下ろして「ワイヤレスキャスト」を探す。
  \item お気に入りの写真を一枚テレビにキャストし、家族に見せましょう。
  \item スマートホームデバイスがあれば、スマホAppで「外出モード」を作成してみる：ワンタップで照明オフ、エアコンオフ、ロボット掃除機起動。
  \item 第十八章：データバックアップと移行
\end{enumerate}

\end{tcbpractice}

\begin{tcbinfo}{事例1：バッテリー持ち改善の実戦}
\begin{itemize}
  \item \textbf{シーン}: ユーザーの Xiaomi 14 バッテリー持ちが悪く、半日で切れる。
  \item \textbf{手順}: (1) 設定 バッテリー 消費ランキング確認 (2) 異常なバックグラウンド App を特定 (3) 5G を 4G に切替 (4) バックグラウンド制限 (5) 省電力モード有効化 (6) 充電時にケースを外して放熱。
  \item \textbf{期待される結果}: バッテリー持ちが30-50\%延長、バッテリー劣化度データが読み取り可能。
\end{itemize}

\end{tcbinfo}
